Examining the three case studies alongside the eight other successful redistricting initiatives in the historical record, we identify five design elements common to successful initiatives and three political barriers that have defeated initiatives in states where redistricting reform was attempted.

\subsection{Five Common Design Elements}

\textbf{(1) Independent commission structure with conflict-of-interest rules.} Every successful initiative created an independent body (not the legislature) to draw maps, with explicit rules excluding elected officials, their families, political consultants, major donors, and lobbyists from commission membership. The conflict-of-interest rules must be specific enough to prevent evasion (e.g., a former elected official who returns to civilian life within one or two years before the redistricting cycle should be excluded) and strong enough to withstand constitutional challenge.

\textbf{(2) Explicit partisan-balance requirements.} Every successful commission design specifies the partisan composition of the commission (or, in Michigan's case, an explicitly partisan-balanced selection pool). Pure neutrality---appointing commissioners without reference to party---is not sufficient; it creates selection processes that incumbents can control. Partisan balance is a structural guarantee of partisan neutrality: no single party can dominate the commission's decisions.

\textbf{(3) Transparent public engagement mandates.} Every successful initiative required public hearings before, during, and after the drawing of maps, with specific requirements (minimum number of hearings, geographic distribution, public comment periods). Transparency serves two purposes: it builds public legitimacy for the final maps, and it creates a record that courts can evaluate when maps are challenged.

\textbf{(4) Enumerated criteria with priority ordering.} Successful initiatives list the criteria a map must satisfy and specify their priority when they conflict. California's six-criteria hierarchy (federal requirements first, VRA compliance second, compactness and contiguity third and fourth, political subdivision respect fifth) is the most widely cited model. The priority ordering is as important as the criteria themselves: without it, mapmakers can select which criteria to optimize by choosing which criteria to characterize as ``primary.''

\textbf{(5) Constitutional entrenchment.} Initiatives that amended the state constitution survived legislative counter-measures more reliably than those enacted as statutes. California's Proposition 11 amended the constitution; when the legislature passed a bill attempting to limit CCRC authority in 2012, the amendment's constitutional status provided a strong defensive posture. Michigan's Proposal 2 is a constitutional amendment; Colorado's Amendments Y and Z are also constitutional amendments. All three were designed to be resistant to simple legislative reversal.

\subsection{Three Political Barriers}

\textbf{(1) Incumbent spending.} The single strongest predictor of initiative failure is high incumbent spending against the initiative. California's Proposition 11 nearly failed due to \$1M+ in opposition spending; Arizona's Proposition 203 (2000) failed after \$1.5M in opposition from the legislature. Reformers must anticipate this spending and either out-fundraise or design initiatives that are difficult to frame as partisan threats.

\textbf{(2) Confusing ballot language.} Redistricting is a complex topic, and initiatives are often written by lawyers who produce ballot language that is accurate but impenetrable. Several failed initiatives (Ohio 2015, Missouri 2020) suffered from public confusion about what the initiative actually did. Short, clear ballot language---``establishes an independent commission to draw congressional districts''---outperforms technically precise but lengthy descriptions.

\textbf{(3) Coalition fracture over algorithmic requirements.} The most relevant barrier for the model initiative in Section~4 is that algorithmic redistricting can fracture the reform coalition. Good-government groups who support commission reform may be skeptical of algorithmic mandates; communities of color may fear that algorithms will override VRA requirements; partisan reformers may worry that the algorithm will systematically disadvantage their party. Each of these concerns is addressable by design (the model language in Section~4 addresses all three), but the framing of the initiative must be managed carefully.
